\documentclass[11pt]{article}
\usepackage[margin=1in]{geometry}
\usepackage{setspace}
\usepackage{hyperref}
\hypersetup{colorlinks=true, urlcolor=blue}
\setlength{\parindent}{0pt}
\setlength{\parskip}{0.9em}
\onehalfspacing
\pagestyle{empty}

\title{\vspace{-1.2em}\textbf{Teaching Statement}}
\author{HSU Ke Cheng\thanks{PhD Candidate in Economics, Hong Kong University of Science and Technology, email: \href{mailto:khsu@connect.ust.hk}{khsu@connect.ust.hk}}}
\date{\vspace{-2.5em}}

\begin{document}
\maketitle

Teaching is most fulfilling when I see a student move from observing a pattern to understanding the mechanism behind it. Economic models make that shift possible: they help students separate what they see from what is causing it. Consider a new restaurant on campus. Before students have tasted the food, its long line appears to indicate quality. Yet the line is only a signal. As each person observes it and joins, the line can continue to grow even when the food is bad. Or consider why a government might issue consumption vouchers rather than cash. Because a voucher must be spent, more of the transfer enters current demand, producing a larger multiplier. I begin with familiar situations like these and establish the economic intuition before introducing a formula, because the main hurdle is often the meaning rather than the algebra. I then ask students to work through the reasoning themselves and to identify the point at which it stops being clear. One student's question often reveals a confusion shared by several others, benefiting the whole class while showing me where my explanation can be simpler. Finding that clearer explanation is the part of teaching I enjoy most.

My teaching experience at the Hong Kong University of Science and Technology has developed this approach. From 2020 to 2025, I served as a teaching assistant for undergraduate macroeconomics, industrial organization, and the economics of uncertainty and information, leading weekly tutorials, holding office hours, and grading. In tutorials and office hours, I learned to identify the precise point at which a student's reasoning became unclear and rebuild the argument from there. In PhD microeconomic theory, taken by doctoral students across the Business School, I primarily led tutorials and answered questions, sometimes approaching the same question in several ways until the idea became clear. On an evaluation scale where 4 means good and 5 means extremely satisfied, my average across these appointments is 4.48. My two evaluations for PhD microeconomic theory were 4.85 and 4.83. In 2026, I lectured in the mathematics camp for the MSc in Economics. Most participants did not have the mathematical background assumed by the program, so I built the intuition and notation step by step. My teaching received a score of 4.5 out of 5, and student comments included ``Very nice and patient!'' and ``Excellent teaching.'' From 2025 to 2026, I also served as Teaching Assistant Coordinator for the Business School. I helped postgraduate students design teaching materials, evaluated their mock teaching, and discussed where students might lose the thread and how those steps could be explained more clearly.

As the instructor of my own courses, I would bring the same emphasis on intuition, active reasoning, and clear explanation. I am prepared to teach microeconomic theory from the undergraduate level through the PhD level; industrial organization, information economics, structural estimation, and mathematics for economists from the undergraduate level through the master's level; and empirical industrial organization, macroeconomics, and econometrics at the undergraduate level. I am also eager to teach students how to manage a research project with AI. An AI system can produce an answer that appears polished but is unreliable when the problem has not been defined precisely. I would therefore teach students to specify the problem and decide in advance what evidence will show that the task is complete. I would also teach them a simple workflow: report their progress, maintain clear code in R or Python, and use Git to track versions. These habits help students catch mistakes and verify the work produced with AI.

\end{document}
